\documentclass[11pt,a4paper]{article}

% --- Packages ---
\usepackage[margin=.75in]{geometry}
\usepackage[utf8]{inputenc}
\usepackage[T1]{fontenc}
\usepackage{lmodern} % Fix for font not found error
\usepackage{titlesec}
\usepackage{enumitem}
\usepackage{hyperref}
\usepackage{xcolor}
\usepackage{fancyhdr}
\usepackage{comment}
\usepackage{cleveref}
\usepackage[bottom]{footmisc}
\usepackage{graphicx}
\usepackage{longtable}
\usepackage{booktabs}
\usepackage{array}
\usepackage{calc}
\usepackage{etoolbox}

% Standard Pandoc setup
\providecommand{\tightlist}{%
  \setlength{\itemsep}{0pt}\setlength{\parskip}{0pt}}

\crefformat{section}{\S#2#1#3}
\crefformat{subsection}{\S#2#1#3}
\crefformat{subsubsection}{\S#2#1#3}
\crefformat{figure}{#2Figure~#1#3}
\crefformat{footnote}{#2\footnotemark[#1]#3}

\linespread{0.9} % Riduce lo spazio tra le linee

% --- Colors ---
\definecolor{primary}{RGB}{0, 0, 0}
\definecolor{links}{HTML}{0000EE}
\definecolor{blue}{HTML}{26428B}

\hypersetup{
    colorlinks=true,
    linkcolor=links,
    urlcolor=links,
    citecolor=links,
}

% --- Section Formatting ---
\titleformat{\section}
  {\normalfont\Large\bfseries\color{blue}}{\thesection}{1em}{}
\titleformat{\subsection}
{\normalfont\large\bfseries\color{blue}}{\thesubsection}{1em}{}
\titleformat{\subsubsection}
{\normalfont\bfseries\color{blue}}{\thesubsubsection}{1em}{}

% --- Header and Footer ---
\pagestyle{fancy}
\fancyhead[L]{\small Giuseppe Capasso}
\fancyhead[R]{\small intro-to-linux}
\fancyfoot[C]{\thepage}
\renewcommand{\headrulewidth}{0.4pt}

% --- Custom Title Command ---
\newcommand{\proposaltitle}[1]{
    \begin{center}
        \vspace*{0.1cm}
        {\LARGE \bfseries \color{blue} #1} \\[1.5ex]
        {\large \textbf{Giuseppe Capasso}} \\[1ex]
                \small{
            \vspace{0.1cm}
            \textbf{Materia:} intro-to-linux\\
        }
            \end{center}
    \vspace{0.3cm}
    \hrule height 0.5pt
    \vspace{0.5cm}
}

\begin{document}

\proposaltitle{Blocca tutto in entrata}

Dispensa Tecnica: Servizi, SSH e Sicurezza di Base Modulo: Linux System
Administration - Lezione 5 Obiettivo: Gestire un server ``Headless''
(senza monitor), configurare l'accesso remoto sicuro, gestire i servizi
di sistema (Demoni) e configurare il Firewall.

\begin{enumerate}
\def\labelenumi{\arabic{enumi}.}
\item
  Benvenuti nel ``Vero'' Linux (Ubuntu Server) Fino ad ora abbiamo usato
  una versione Desktop con interfaccia grafica. Tuttavia, nel mondo
  reale, i server Linux (Cloud, Datacenter) sono Headless: non hanno
  monitor, mouse o finestre. Tutto si gestisce tramite riga di comando
  da remoto. Perché? • Risorse: L'interfaccia grafica consuma RAM e CPU
  inutilmente. • Sicurezza: Meno software installato significa meno bug
  potenziali. • Automazione: La CLI è facile da scriptare, la GUI no.
\item
  Secure Shell (SSH) Poiché non possiamo sederci fisicamente davanti al
  server, usiamo SSH. È un protocollo che crea un ``tunnel''
  crittografato tra il tuo computer (Client) e il server.
\end{enumerate}

2.1 Autenticazione: Password vs Chiavi SSH supporta due metodi
principali di accesso: 1. Password: Semplice, ma vulnerabile agli
attacchi ``Brute Force'' (hacker che provano milioni di password al
secondo). 2. Chiavi Asimmetriche (Key-Based): Metodo standard
professionale. Si basa su una coppia di file crittografici. • Chiave
Privata (Private Key): Risiede sul tuo PC (Client). È la tua ``chiave di
casa''. Non va mai condivisa. • Chiave Pubblica (Public Key): Risiede
sul Server. È come la ``serratura''. Puoi darla a chiunque. Il server ti
fa entrare solo se la tua chiave privata ``apre'' la chiave pubblica che
hai installato lì.

2.2 Comandi SSH Fondamentali Comando (sul tuo PC) Descrizione ssh
user@192.168.1.50 Si collega al server specificato. ssh-keygen -t
ed25519 Genera una nuova coppia di chiavi sul tuo PC. ssh-copy-id
user@IP Copia la tua chiave pubblica sul server remoto. Esporta in Fogli

\begin{enumerate}
\def\labelenumi{\arabic{enumi}.}
\setcounter{enumi}{2}
\tightlist
\item
  Gestione dei Servizi (Systemd) Un Servizio (o Demone) è un programma
  che gira in background, senza interfaccia utente, in attesa di
  richieste (es. Web Server, Database, SSH). In Linux moderni, il
  ``capo'' che gestisce questi servizi si chiama Systemd. Il comando per
  controllarlo è systemctl.
\end{enumerate}

3.1 Ciclo di vita di un servizio Azione

Comando Significato sudo systemctl start Avviare apache2 Accende il
servizio ora (fino al prossimo riavvio). sudo systemctl stop Fermare
apache2 Spegne il servizio immediatamente. sudo systemctl restart
Riavviare apache2 Stop + Start (utile quando cambi le configurazioni).
sudo systemctl status Stato Ti dice se è ``Active (running)'' o
``Failed''. apache2 sudo systemctl enable Fondamentale: Dice al servizio
di partire Abilitare apache2 automaticamente all'avvio del server.
Esporta in Fogli

\begin{enumerate}
\def\labelenumi{\arabic{enumi}.}
\setcounter{enumi}{3}
\tightlist
\item
  Sicurezza di Rete: Firewall (UFW) Linux ha un firewall integrato
  potentissimo (Netfilter/Iptables), ma difficile da configurare. Noi
  useremo UFW (Uncomplicated Firewall), un'interfaccia semplificata.
\end{enumerate}

4.1 Principio del ``Default Deny'' La regola d'oro della sicurezza è:
Blocca tutto ciò che entra, tranne quello che serve esplicitamente.

4.2 Configurazione UFW 1. Imposta le policy di base: Bash sudo ufw
default deny incoming

\section{Blocca tutto in entrata}\label{blocca-tutto-in-entrata}

sudo ufw default allow outgoing

\section{Permetti tutto in uscita}\label{permetti-tutto-in-uscita}

\begin{enumerate}
\def\labelenumi{\arabic{enumi}.}
\setcounter{enumi}{1}
\tightlist
\item
  Apri le porte necessarie (Whitelist): Bash sudo ufw allow ssh
\end{enumerate}

\section{(Porta 22) - FONDAMENTALE FARLO PRIMA DI
ATTIVARE!}\label{porta-22---fondamentale-farlo-prima-di-attivare}

sudo ufw allow http

\section{(Porta 80)}\label{porta-80}

\begin{enumerate}
\def\labelenumi{\arabic{enumi}.}
\setcounter{enumi}{2}
\tightlist
\item
  Attiva il firewall: Bash sudo ufw enable
\end{enumerate}

ATTENZIONE: Se attivi UFW senza aver prima dato allow ssh, verrai
buttato fuori dal server e non potrai più rientrare da remoto!

\begin{enumerate}
\def\labelenumi{\arabic{enumi}.}
\setcounter{enumi}{4}
\tightlist
\item
  Gestione dei Log (Troubleshooting) Quando qualcosa non funziona, Linux
  scrive il motivo in un file di testo. I log si trovano quasi sempre in
  /var/log.
\end{enumerate}

5.1 File di Log Critici • /var/log/syslog (o messages): Il diario
generale del sistema. • /var/log/auth.log: Log di Sicurezza. Registra
chi entra, chi esce e chi sbaglia la password (SSH, sudo). •
/var/log/apache2/access.log: Chi sta visitando il tuo sito web. •
/var/log/apache2/error.log: Perché il tuo sito web non funziona.

5.2 Analisi in Tempo Reale Non aprire i log con un editor di testo! Sono
enormi. Usa tail. Bash \# Segue il file in diretta mentre vengono
scritte nuove righe tail -f /var/log/auth.log

(Premi CTRL+C per uscire)

Laboratorio Pratico: Server Hardening Scenario: Hai appena installato
una nuova macchina Ubuntu Server. Devi metterla in sicurezza,
configurare un web server e disabilitare l'accesso via password.

Parte 1: Setup e Accesso SSH 1. Avvia la VM Ubuntu Server e fai login.
2. Trova il tuo IP: ip addr show. 3. DA ORA IN POI, lavora dal terminale
del tuo PC reale (Host): • Collegati: ssh studente@Tuo\_IP. • Accetta il
fingerprint (scrivi yes). • Inserisci la password.

Parte 2: Configurazione Chiavi (Key-Based Auth) 1. Sul tuo PC Host (apri
un altro terminale): • Genera le chiavi: ssh-keygen -t ed25519 (Premi
Invio a tutte le domande). • Copia la chiave sulla VM: ssh-copy-id
studente@Tuo\_IP. 2. Prova: Prova a collegarti di nuovo via SSH. Se non
ti chiede la password, hai fatto tutto correttamente.

Parte 3: Blindare SSH Ora che le chiavi funzionano, disabilitiamo le
password per evitare attacchi hacker. 1. Sulla VM, edita la
configurazione SSH: Bash sudo nano /etc/ssh/sshd\_config

\begin{enumerate}
\def\labelenumi{\arabic{enumi}.}
\setcounter{enumi}{1}
\tightlist
\item
  Trova la riga \#PasswordAuthentication yes, togli il commento (\#) e
  cambiala in: PasswordAuthentication no
\item
  Salva ed esci (CTRL+O, CTRL+X).
\item
  Riavvia il servizio per applicare le modifiche: Bash sudo systemctl
  restart ssh
\end{enumerate}

Parte 4: Web Server e Servizi 1. Installa Nginx: sudo apt update \&\&
sudo apt install nginx. 2. Verifica che sia attivo: systemctl status
nginx. 3. Controlla se è abilitato al boot: systemctl is-enabled nginx.

Parte 5: Configurazione Firewall 1. Imposta le regole: Bash sudo ufw
default deny incoming sudo ufw allow ssh sudo ufw allow http

\begin{enumerate}
\def\labelenumi{\arabic{enumi}.}
\setcounter{enumi}{1}
\tightlist
\item
  Attivalo: sudo ufw enable.
\item
  Verifica lo stato: sudo ufw status.
\end{enumerate}

Parte 6: Analisi Intrusioni (Log) 1. Tieni aperto il collegamento SSH.
2. Prova ad aprire un nuovo terminale dal tuo PC e prova a collegarti
come un utente inesistente: ssh hacker@Tuo\_IP (Fallirà). 3. Torna al
terminale SSH connesso e cerca la traccia del tentativo nel log: Bash
sudo grep ``Invalid user'' /var/log/auth.log

Dovresti vedere il tentativo fallito dell'utente ``hacker''.



\end{document}
